Los experimentos fueron realizados en un computador personal Legion 5 15ITH6 utilizando sistema operativo Linux (WSL sobre Windows mediante Ubuntu 22.04 LTS), con procesador 11th Gen Intel(R) Core(TM) i5-11400H @ 2.70GHz (2.69 GHz) y 8 GB de memoria RAM. El código fue compilado utilizando \texttt{g++} con estándar \texttt{C++11}.

El programa mide automáticamente el tiempo de ejecución mediante la librería \texttt{chrono} y el uso de memoria mediante \texttt{getrusage}. La ejecución de pruebas se automatizó mediante un \texttt{Makefile}, generando resultados reproducibles.

Los datasets utilizados corresponden a instancias sintéticas del problema, variando el número de animes \(n\) en valores: 5, 10, 15, 20, 30, 50, 80, 120 y 200. Para cada tamaño se generaron múltiples casos de prueba, permitiendo calcular promedios y analizar el comportamiento de los algoritmos.

Los resultados se almacenan automáticamente en archivos dentro de \texttt{code/implementation/data/measurements/}, y posteriormente se procesan mediante scripts en Python para generar gráficos en \texttt{data/plots/}.